% Target ~1900 words. Rubric: System Design + System Evaluation (25 of 50 marks).
% This is the section that decides the grade. Every [FILL] below needs real data.

\section{Results: prototypes and evaluation findings}

\subsection{Response to Assignment 2 feedback}
\label{sec:a2response}

Feedback on Assignment 2 was received on 2 October 2026. The summary comment was that the report was
well structured but that much of the content rested on assumptions, and that mid-fidelity wireframes
were lacking for important features and user flows. Table~\ref{tab:a2response} records each point of
feedback and what was changed in response.

\begin{table}[htbp]
\centering
\caption{Assignment 2 feedback and the response made in this iteration.}
\label{tab:a2response}
\small
\begin{tabularx}{\textwidth}{@{}lXX@{}}
\toprule
\textbf{\#} & \textbf{Feedback} & \textbf{Our response} \\
\midrule
1 & Content rested on assumptions rather than evidence &
Claims about usability and motivation are now tested with ten participants (Section~\ref{sec:findings}). \\
2 & Wireframes lacking for important features and flows &
The prototype now covers the whole journey in \VERIFY{number} linked screens, including flows that A2 only described: voice coordination, moderation, and the community view. \\
3 & Unclear how heuristics were applied to improve the design &
Section~\ref{sec:changes} records each change with the heuristic or finding that motivated it. \\
4 & The host role was undefined (``is this person just there monitoring?'') &
The host now has defined powers and a screen: reading the coverage radar and assigning stations. \\
5 & Collaboration was thin (``is it just taking photos?'') &
\FILL{what changed --- shared radar state, assignments, quick callouts} \\
6 & Session governance unspecified (which places, who decides, time limits, clashes) &
\FILL{what changed --- quorum, geofence, the two-groups-one-landmark flow} \\
7 & Weak user motivation and engagement &
\FILL{what changed --- profile, reputation, condition sets} \\
8 & Premise under-evidenced; references too few or arbitrary &
\FILL{what changed --- citations added in Sections 1, 2 and 4} \\
\bottomrule
\end{tabularx}
\end{table}

\VERIFY{Rows 4--8 come from the marker's annotated comments, transcribed in
\texttt{docs/A2-FEEDBACK.md}. Quote each comment briefly in the Feedback column rather than
paraphrasing, and make sure every row has a real response. Appendix~\ref{app:a2feedback} holds the
full set.}

\subsection{Prototype evolution}

\subsubsection{Initial prototypes (Assignment 2)}

The handout requires images of all initial prototypes. Figures~\ref{fig:sketch1}
to~\ref{fig:lofi} reproduce the paper sketches, storyboard and low-fidelity wireframes produced
during Assignment 2.

\begin{figure}[htbp]
\centering
\safeincludegraphics[width=0.78\textwidth]{fig1-scenario1-join-session.png}
\caption{Paper sketch, scenario 1: joining a session by code, seeing the shared start time,
capturing, receiving the model and sharing it.}
\label{fig:sketch1}
\end{figure}

\begin{figure}[htbp]
\centering
\safeincludegraphics[width=0.78\textwidth]{fig2-scenario2-qr-recruitment.png}
\caption{Scenario 2: recruiting passers-by on the spot with a QR code on a poster.}
\label{fig:sketch2}
\end{figure}

\begin{figure}[htbp]
\centering
\safeincludegraphics[width=0.78\textwidth]{fig3-scenario3-create-session.png}
\caption{Scenario 3: creating a session with date, time, location and duration, then sharing the
code.}
\label{fig:sketch3}
\end{figure}

\begin{figure}[htbp]
\centering
\safeincludegraphics[width=0.46\textwidth]{fig4a-home-and-create-sketch.png}\hfill
\safeincludegraphics[width=0.46\textwidth]{fig4b-camera-and-share-sketch.png}
\caption{Paper sketches of the four main screens: home and create session (left); camera and share
(right).}
\label{fig:sketch4}
\end{figure}

\begin{figure}[htbp]
\centering
\safeincludegraphics[width=0.9\textwidth]{fig5-storyboard.png}
\caption{Scenario storyboard: three strangers discover ChronoMap through a QR poster, form a
session, coordinate angles through coverage prompts, and receive contribution scores.}
\label{fig:storyboard}
\end{figure}

\begin{figure}[htbp]
\centering
\safeincludegraphics[width=0.95\textwidth]{fig6-lofi-figma-wireframes.png}
\caption{Low-fidelity Figma wireframes of the four core screens: discovery and lobby, coverage
guide, capture window, and result and sharing.}
\label{fig:lofi}
\end{figure}

In these early rounds a session was simply a shared time window and each participant pressed the
shutter freely. Reviewing the sketches exposed that nothing told people where to stand and nothing
showed how much of the site had been covered, which is the problem the high-fidelity design sets out
to solve.

\subsubsection{High-fidelity prototype, version 1}

Version 1 translated the low-fidelity wireframes into a complete, linked journey. It is the build
that was put in front of participants.

\FILL{Describe version 1: the screens, the journey, and what changed from the low-fidelity round
and why. Reference the design decision log. This is the ``several iterations with motivated changes'' the rubric asks for, so each change
needs a stated reason rather than a description alone.}

\begin{table}[htbp]
\centering
\caption{Changes between prototype rounds, and what motivated each.}
\label{tab:iterations}
\small
\begin{tabularx}{\textwidth}{@{}llXX@{}}
\toprule
\textbf{Round} & \textbf{Change} & \textbf{Motivation} & \textbf{Source} \\
\midrule
Sketch $\rightarrow$ storyboard & \FILL{} & \FILL{} & Sketch review \\
Storyboard $\rightarrow$ low-fi & \FILL{} & \FILL{} & Storyboard walkthrough \\
Low-fi $\rightarrow$ hi-fi v1 & \FILL{} & \FILL{} & A2 feedback \\
Within hi-fi v1 & \FILL{} & \FILL{} & Heuristic review \\
Hi-fi v1 $\rightarrow$ v2 & \FILL{each change made after the evaluation} & \FILL{finding ID} & Evaluation \\
\bottomrule
\end{tabularx}
\end{table}

\subsection{Features and the social behaviours they support}

\FILL{Walk the main journey and explain how each feature supports social behaviour. The rubric
rewards explaining \emph{how} features support social behaviour, not just listing them. Cover at
minimum: the quorum gate, station assignment, the shared coverage radar, quick callouts, moderation,
and the community view.}

\subsection{Evaluation findings}
\label{sec:findings}

\subsubsection{Task performance}

\begin{table}[htbp]
\centering
\caption{Task completion by participant (think-aloud group).}
\label{tab:taskperf}
\small
\begin{tabular}{@{}llccc@{}}
\toprule
\textbf{Task} & \textbf{Completed} & \textbf{Assists} & \textbf{Median time} & \textbf{Notes} \\
\midrule
T1 & \FILL{} & \FILL{} & \FILL{} & \FILL{} \\
T2 & \FILL{} & \FILL{} & \FILL{} & \FILL{} \\
T3 & \FILL{} & \FILL{} & \FILL{} & \FILL{} \\
T4 & \FILL{} & \FILL{} & \FILL{} & \FILL{} \\
T5 & \FILL{} & \FILL{} & \FILL{} & \FILL{} \\
\bottomrule
\end{tabular}
\end{table}

\subsubsection{Questionnaire results}

\FILL{SUS mean, standard deviation and range, against the published average of 68~\cite{brooke1996}.
Say which items scored lowest and what that implies.}

\FILL{The six custom items: willingness to join strangers, comfort with voice, radar clarity, trust
in privacy controls, motivation, purpose comprehension. Report means and note any item where the
spread matters more than the average.}

\FILL{Themes from the two open questions, with a short representative quote for each.}

\subsubsection{Consolidated findings}

\begin{table}[htbp]
\centering
\caption{Usability problems found, with severity and the decision taken.}
\label{tab:findings}
\small
\begin{tabularx}{\textwidth}{@{}llXlll@{}}
\toprule
\textbf{ID} & \textbf{Screen} & \textbf{Problem} & \textbf{P's} & \textbf{Sev.} & \textbf{Decision} \\
\midrule
F1 & \FILL{} & \FILL{} & \FILL{} & \FILL{} & \FILL{} \\
F2 & \FILL{} & \FILL{} & \FILL{} & \FILL{} & \FILL{} \\
F3 & \FILL{} & \FILL{} & \FILL{} & \FILL{} & \FILL{} \\
F4 & \FILL{} & \FILL{} & \FILL{} & \FILL{} & \FILL{} \\
F5 & \FILL{} & \FILL{} & \FILL{} & \FILL{} & \FILL{} \\
\bottomrule
\end{tabularx}
\end{table}

Severity follows Nielsen's scale~\cite{nielsen1994}, from 0 (not a problem) to 4 (prevents task
completion). A finding marked \emph{defer} or \emph{reject} carries a stated reason; an honest
deferral is more defensible than a silently ignored problem.

\subsection{Changes made in response to the evaluation (version 2)}
\label{sec:changes}

Version 2 is the prototype submitted with this report. \FILL{For each change: what was altered,
which finding drove it, and a before/after figure pair.
This subsection is what the handout means by ``make improvements to your high-fidelity prototype
based on lessons learned'', and it is the difference between a described evaluation and a used one.}

% Template for each before/after pair - duplicate as needed.
% \begin{figure}[htbp]
% \centering
% \includegraphics[width=0.45\textwidth]{before-XX.png}\hfill
% \includegraphics[width=0.45\textwidth]{after-XX.png}
% \caption{Finding F\#: before (left) and after (right). \FILL{what changed and why}}
% \end{figure}
